\documentclass[11pt,a4paper]{article}
\usepackage[margin=1in]{geometry}
\usepackage{amsmath,amssymb,mathtools}
\usepackage{bm}
\usepackage{microtype}
\usepackage{enumitem}

\newcommand{\R}{\mathbb{R}}
\newcommand{\E}{\mathbb{E}}
\newcommand{\softmax}{\mathrm{softmax}}
\newcommand{\diag}{\mathrm{diag}}

\title{PropertyScanner: Cross-Attention Comp Fusion for Market Valuation}
\author{}
\date{}

\begin{document}
\maketitle

\begin{abstract}
This note documents the valuation model and evaluation protocol used in
PropertyScanner. I focus on the mathematics that connects market reasoning with
modern representation learning: comparable retrieval, time-adjusted anchors,
cross-attention fusion, quantile regression, and conformal calibration. The
goal is not only prediction but a defensible interval around fair value.
Throughout, I adopt notation that keeps the economics visible, because in my
experience in financial modeling, a model without an anchor to market structure
is a model that drifts.
\end{abstract}

\section{Setup and Notation}
Each listing $i$ has observed price $y_i > 0$ and three modalities:
tabular features $x_i \in \R^{d_x}$, text embedding $t_i \in \R^{d_t}$, and
optional image embedding $v_i \in \R^{d_v}$.
In the current implementation, $d_x = 11$, $d_t = 384$, $d_v = 512$.
The tabular vector includes location, size, floor, sentiment, and a
price-per-sqm proxy; text includes title, description, and optionally
vision-language descriptions.

Given a target listing $0$, we retrieve $K$ comparables
$\mathcal{C} = \{1,\ldots,K\}$ with prices $y_j$ and features
$(x_j,t_j,v_j)$.
Let $r$ denote the region (city) of the target and let $m(i)$ denote the month
of listing $i$.
The overall pipeline is:
\[
\text{retrieve comps} \rightarrow \text{time-adjust prices} \rightarrow
\text{cross-attention fusion} \rightarrow \text{quantile outputs} \rightarrow
\text{calibration} \rightarrow \text{evaluation}.
\]

\section{Comparable Retrieval}
Semantic retrieval uses dense text embeddings with a FAISS index, followed by
structural filters. Let $e(\cdot)$ be the text encoder. A query vector
$q = e(\text{title}_0 \oplus \text{desc}_0 \oplus \text{vlm}_0)$ is compared to
candidate vectors by inner product (normalized cosine similarity).
The candidate pool is then filtered by distance, property type, and size:
\begin{align}
\text{dist}(j,0) &\leq R, \\
\text{type}(j) &= \text{type}(0), \\
1-\rho \leq \frac{\text{sqm}_j}{\text{sqm}_0} &\leq 1+\rho.
\end{align}
This two-stage filter reflects a key valuation principle: similarity must be
semantic and structural.

\section{Hedonic Time Adjustment}
Comparable prices are time-adjusted to the valuation date using a hedonic index.
The index is estimated by a monthly regression
\begin{equation}
\log P = \beta^\top f + \gamma_{\text{nh}} + \varepsilon,
\end{equation}
where $f$ contains size, bedrooms, bathrooms, elevator, and floor, and
$\gamma_{\text{nh}}$ is a neighborhood fixed effect.
Let $I_r(m)$ be the quality-adjusted price-per-sqm index for region $r$ at month
$m$. For comp $j$ observed at month $m(j)$ and valued at month $m(0)$,
\begin{equation}
y_j^\ast = y_j \cdot \frac{I_r(m(0))}{I_r(m(j))}.
\end{equation}
This prevents leakage from temporal price drift and gives the model a
current-market anchor.

\section{Multimodal Encoding and Fusion}
Each modality is projected into a shared hidden dimension $d$:
\begin{align}
h_x &= \phi_x(W_x x + b_x), \\
h_t &= W_t t, \\
h_v &= W_v v, \\
h &= \phi_f([h_x;h_t;h_v]),
\end{align}
with $d=64$ in the deployed model. The projection layers are shallow by design:
they align features without erasing structure, a common failure mode in
over-parameterized tabular fusion.

For the target and its comps, define $h_0 \in \R^d$ and
$H \in \R^{K \times d}$.
Cross-attention uses the target as query and comps as keys and values:
\begin{align}
A &= \softmax\left(\frac{h_0 H^\top}{\sqrt{d}}\right) \in \R^{1 \times K}, \\
z &= A H \in \R^{1 \times d}, \\
r &= g(h_0 + z),
\end{align}
where $g(\cdot)$ is a small post-attention MLP.
The attention weights $A$ also serve as interpretable comp weights.

\section{Anchored Residual Prediction}
From a finance perspective, the safest predictive strategy is to anchor to
market evidence and learn residual adjustments. We form an attention-weighted
price anchor
\begin{equation}
a = \sum_{j=1}^K A_j y_j^\ast,
\end{equation}
and a scale term
\begin{equation}
s = \mathrm{std}(y_1^\ast,\ldots,y_K^\ast).
\end{equation}
For quantile $q \in \{0.1,0.5,0.9\}$, the model predicts a residual
${u}_q = w_q^\top r + b_q$, and outputs
\begin{equation}
\hat{y}_q = a + s \cdot {u}_q.
\end{equation}
This simple decomposition keeps the model honest: it can only move predictions
relative to what the market comps support, and the residual scale adapts to
local dispersion.

\section{Training Objective}
We use quantile regression with the pinball loss. For a single target with
prediction $\hat{y}_q$:
\begin{equation}
\ell_q(y,\hat{y}_q) =
\max\left(q(y-\hat{y}_q), (q-1)(y-\hat{y}_q)\right).
\end{equation}
The training loss is the mean across the three quantiles:
\begin{equation}
\mathcal{L} = \frac{1}{3} \sum_{q \in \{0.1,0.5,0.9\}} \ell_q(y,\hat{y}_q).
\end{equation}
Only price quantiles are optimized in the current training loop; other heads
are deployed but not yet trained on explicit labels.
Optimization uses AdamW with gradient clipping and early stopping.
The dataset samples comps using a geo-radius and size filter, and can be
evaluated using a standard holdout split or K-fold cross validation.

\section{Uncertainty and Calibration}
The raw uncertainty proxy is the relative half-width:
\begin{equation}
u = \frac{\hat{y}_{0.9} - \hat{y}_{0.1}}{2\hat{y}_{0.5}}.
\end{equation}
Intervals are then calibrated by adaptive conformal prediction per horizon.
Given residuals
\begin{align}
e^- &= \max(0, \hat{y}_{0.1} - y), \\
e^+ &= \max(0, y - \hat{y}_{0.9}),
\end{align}
we compute quantiles $c^- = Q_{1-\alpha}(e^-)$ and
$c^+ = Q_{1-\alpha}(e^+)$ on a rolling window, and widen the interval:
\begin{equation}
\tilde{y}_{0.1} = \hat{y}_{0.1} - c^-, \quad
\tilde{y}_{0.9} = \hat{y}_{0.9} + c^+.
\end{equation}
This asymmetric calibration matches market intuition: downside and upside errors
behave differently in illiquid markets.

\section{Evaluation Protocol}
Evaluation is organized around two axes: point accuracy and interval quality.
Point accuracy uses the median prediction $\hat{y}_{0.5}$:
\begin{align}
\mathrm{MAE} &= \frac{1}{N}\sum_{i=1}^N |y_i - \hat{y}_{i,0.5}|, \\
\mathrm{MAPE} &= \frac{100}{N}\sum_{i=1}^N
\frac{|y_i - \hat{y}_{i,0.5}|}{\max(y_i,1)}, \\
\mathrm{MedAE} &= \mathrm{median}_i |y_i - \hat{y}_{i,0.5}|.
\end{align}
Interval quality is evaluated by empirical coverage and width:
\begin{align}
\mathrm{Coverage} &= \frac{1}{N} \sum_{i=1}^N
\mathbf{1}\{\tilde{y}_{i,0.1} \le y_i \le \tilde{y}_{i,0.9}\}, \\
\mathrm{Width} &= \frac{1}{N} \sum_{i=1}^N
(\tilde{y}_{i,0.9} - \tilde{y}_{i,0.1}).
\end{align}
These metrics align with the deployment goal: the model must be accurate enough
to rank opportunities, and calibrated enough to communicate risk.

\section{Closing Perspective}
The model is intentionally modest in parameter count but precise in market
structure. From long experience in financial modeling, the hardest part is not
learning a price function; it is respecting how the market sets prices today,
and how uncertainty should be priced into decisions. Cross-attention delivers a
modern inductive bias for comp reasoning, and conformal calibration turns that
reasoning into intervals that survive regime shifts. That combination is the
core engineering thesis of PropertyScanner.

\end{document}
